\begin{textsample}{POS Dim 6 – rj – Score 7.00 – rj/rj\_em\_2.txt}  \label{ex:f6_pos_272}
PARTE DIVERSIFICADA E EIXOS \textbf{ESTRUTURANTES} De acordo com o Artigo 7º da Resolução SEEDUC Nº 6035, de 28 de janeiro de 2022, que fixa diretrizes para a implantação das \textbf{matrizes} curriculares para a educação básica nas unidades escolares da rede pública : > “ A Parte Diversificada ( Itinerário \textbf{Formativo} ou Núcleo Articulador ) é componente obrigatório do currículo escolar, devendo estar organicamente articulada à Base Nacional Comum Curricular, tornando o currículo um todo significativo e integrado. ” Dentro desse contexto, o Currículo Referencial do Estado do Rio de Janeiro – Ensino Médio pretende fortalecer o protagonismo dos jovens estudantes, possibilitando que escolham o itinerário \textbf{formativo} no qual pretendem aprofundar seus estudos.
Os itinerários \textbf{formativos} representam a parte flexível do Currículo e são estruturados, com base na BNCC, em quatro eixos \textbf{estruturantes} a saber : investigação científica, processos criativos, mediação/intervenção sociocultural e \textbf{empreendedorismo}.
Os itinerários \textbf{formativos} baseados em eixos \textbf{estruturantes} permitem a vivência de situações de aprendizagens práticas, o desenvolvimento de projetos de intervenção na realidade social, com produção de novos conhecimentos.
Cabe ressaltar que, de acordo com o parágrafo único do Artigo 7º, acima mencionado, o planejamento da Parte Diversificada ( Itinerário \textbf{Formativo} ou Núcleo Articulador ) deverá constar no Projeto Político-Pedagógico da escola, priorizando o exercício da autonomia e retratando a identidade da unidade escolar.
Entenda um pouco mais sobre os eixos \textbf{estruturantes} nos textos a seguir.
Investigação Científica A Investigação Científica da realidade por meio de práticas que utilizam os conhecimentos \textbf{estruturantes} das Ciências permite o desenvolvimento das habilidades do pensar e do fazer científico ; Capacita o estudante para compreender e resolver situações cotidianas ; Promove o desenvolvimento local e a melhoria da qualidade de vida da comunidade ; Permite o desenvolvimento e aplicação de produtos desenvolvidos com base científica na realidade cotidiana.
Processos Criativos A idealização e execução de projetos criativos é estabelecida a partir dos conhecimentos da arte, da cultura, da mídia, das ciências e suas aplicações, permitindo o desenvolvimento de habilidades do pensar e do fazer criativo ; Capacita o estudante para expressar -se criativamente e/ou construir soluções inovadoras para a resolução de \textbf{problemas} da sociedade e do mundo do trabalho. \textbf{Mediação} e Intervenção Sociocultural O envolvimento na vida pública, via projetos de mobilização e intervenção sociocultural e \textbf{ambiental}, resulta do conhecimento das questões que afetam a vida dos seres humanos e do planeta.
Esse eixo \textbf{estruturante} permite o desenvolvimento de habilidades de convivência, atuação sociocultural e \textbf{ambiental} ; Capacita o estudante para a \textbf{mediação} de conflitos e proposição de soluções para \textbf{problemas} da comunidade. \textbf{Empreendedorismo} A criação de empreendimentos pessoais ou produtivos articulados ao projeto de vida é fundamental para que ocorra a construção de novos conhecimentos ; Possibilita a percepção do contexto e do mundo do trabalho ; Viabiliza a gestão de iniciativas empreendedoras ; Permite o desenvolvimento de habilidades de autoconhecimento, \textbf{empreendedorismo} e projeto de vida ; Capacita o estudante para estruturar iniciativas empreendedoras que fortaleçam atuação como protagonistas da sua trajetória.
O Ministério de Educação² recomenda que os itinerários \textbf{formativos} integrem os eixos supracitados, integrando -os às bases \textbf{estruturantes} do Currículo.
Dessa forma, os discentes poderão desenvolver habilidades diversificadas e relevantes para uma formação integral.

% matched lemmas: ambiental, empreendedorismo, estruturante, formativo, matriz, mediação, problema
\end{textsample}
